\chapter{原型系统设计与实现}

为将前文提出的多模态安全对齐方法落地为可运行能力，本章面向真实对话服务场景设计并实现原型系统。系统以“安全性优先、可用性不显著退化、工程上可部署可观测”为目标，构建从用户输入到安全输出的端到端链路，并通过典型场景回放展示方法的应用效果。

\section{问题与系统目标}

第四章给出了安全对齐训练方法，第六章将从离线评测角度验证模型性能。本章关注的问题是：如何将训练得到的安全模型封装为可在线服务、可扩展、可维护的系统能力。

围绕该问题，原型系统设定以下三个目标：
\begin{itemize}
    \item \textbf{安全目标：} 对高风险输入具备稳定识别与拦截能力，在边界场景下优先保证输出合规；
    \item \textbf{可用性目标：} 在安全约束下尽可能保持回答有用性，降低过度拒答；
    \item \textbf{工程目标：} 提供统一调用接口与可观测部署形态，支持后续在真实业务环境平滑迁移。
\end{itemize}

据此，系统采用“前置风险判别 + 安全生成 + 事实校验”的分层处理范式，将模型能力映射为明确的工程模块与调用流程。

\section{系统架构与流程}

原型系统采用微服务架构设计，整体可分为三层：接入层、推理层与工具层，如图~\ref{fig:prototype_arch}所示。

\begin{itemize}
    \item \textbf{接入层：} 接收用户多模态输入（文本查询与可选图像），完成格式校验、上下文会话维护与请求编排；
    \item \textbf{推理层：} 包含前置风险分类器与安全对话模型。分类器完成风险类别判定，安全对话模型执行内生风险分析与回复生成；
    \item \textbf{工具层：} 提供检索增强（RAG）能力，在知识密集或高不确定场景中增强事实可信与结果一致性。
\end{itemize}

\begin{figure}[!htb]
    \centering
    \begin{tikzpicture}[
        scale=0.85,
        transform shape,
        node distance=8mm and 12mm,
        >=Latex,
        box/.style={draw, rounded corners, align=center, minimum width=28mm, minimum height=10mm, font=\small},
        layer/.style={draw, dashed, rounded corners, inner sep=4mm},
        arrow/.style={->, thick}
    ]
        % 接入层
        \node[box, fill=blue!8] (api) {API Gateway\\（接入层）};
        \node[box, below=of api, fill=blue!8] (session) {Session管理\\上下文维护};

        % 推理层
        \node[box, below=of session, fill=green!8] (classifier) {前置风险分类器\\（轻量VLM）};
        \node[box, below=of classifier, fill=green!8] (router) {路由决策\\（安全/通用/兜底）};
        \node[box, below left=12mm and 8mm of router, fill=green!15] (general) {通用模型\\回复生成};
        \node[box, below=of router, fill=green!15] (safellm) {安全对话模型\\SafetyCoT+RLVR};
        \node[box, below right=12mm and 8mm of router, fill=green!15] (fallback) {规则兜底\\固定话术};

        % 工具层
        \node[box, below=of safellm, fill=orange!8] (rag) {RAG模块\\检索增强};
        \node[box, right=of rag, fill=orange!8] (retrieval) {检索引擎\\（向量数据库）};

        % 输出
        \node[box, below=of rag, fill=gray!10] (output) {统一输出\\（安全回复/拒绝）};

        % 连线
        \draw[arrow] (api) -- (session);
        \draw[arrow] (session) -- (classifier);
        \draw[arrow] (classifier) -- (router);
        \draw[arrow] (router) -- node[left, font=\scriptsize] {安全} (general);
        \draw[arrow] (router) -- node[right, font=\scriptsize] {风险} (safellm);
        \draw[arrow] (router) -- node[right, font=\scriptsize] {高危} (fallback);
        \draw[arrow] (general) |- (output);
        \draw[arrow] (safellm) -- (rag);
        \draw[arrow] (rag) -- (retrieval);
        \draw[arrow] (rag) -- (output);
        \draw[arrow] (fallback) |- (output);
    \end{tikzpicture}
    \caption{原型系统三层架构与数据流转}
    \label{fig:prototype_arch}
\end{figure}

在端到端处理流程上，系统执行如下步骤：
\begin{enumerate}
    \item 接入层接收请求并完成输入规范化；
    \item 前置分类器输出风险等级标签；
    \item 路由模块按风险等级分流至通用模型、安全模型或兜底策略；
    \item 对知识密集型请求按需触发 RAG，补充外部证据；
    \item 对候选回复执行规则一致性校验并输出最终结果。
\end{enumerate}

\section{核心模块设计}

\subsection{前置风险分类模块}

前置风险分类器采用轻量级大语言模型，通过预缓存提示词模版进行分类任务，依据输出的首个词元对输入查询进行多目标分类：
\textbf{安全正常}、\textbf{一般风险类别}与\textbf{高危违规}。
分类器独立部署，推理延迟控制在 20ms 以内，满足在线服务实时性需求。

分类决策逻辑如下：
\begin{itemize}
    \item \textbf{安全正常：} 直接路由至通用模型回复(放行)，不经过安全对话模型；
    \item \textbf{一般风险：} 路由至安全对话模型，由其进行风险感知与安全生成；
    \item \textbf{高危违规：} 直接触发规则兜底话术，拒绝回答并给出合规提示。
\end{itemize}

\subsection{安全对话生成模块}

安全对话生成模块加载第四章训练得到的 SafetyCoT + RLVR 模型权重。该模块在接收到风险输入后，执行以下流程：

\begin{enumerate}
    \item \textbf{规则检索：} 根据风险类别从规则库检索对应安全原则与实例规则；
    \item \textbf{安全推理：} 模型生成隐式安全反思链，对潜在风险点与用户意图进行分析；
    \item \textbf{约束生成：} 在规则约束下生成兼顾安全性与有用性的回复；
    \item \textbf{后验校验：} 调用规则验证器检查一致性，未通过则触发重写机制。
\end{enumerate}

\subsection{检索增强事实校验模块}

针对知识密集型查询，系统在安全对话链路中集成 Agentic RAG。当模型判断当前问题涉及事实性内容时，由模型自主选择是否调用检索工具：

\begin{enumerate}
    \item 获取模型自主生成的查询，并将查询转化为检索语句，送入向量数据库（Milvus）进行相似度检索；
    \item 获取 Top-K 相关文档片段，作为上下文拼接至生成输入；
    \item 模型在生成回复时显式引用检索内容，降低幻觉风险；
\end{enumerate}

\section{工程实现与部署}

\subsection{统一调用接口}

原型系统对外暴露 RESTful API，核心端点定义如表~\ref{tab:api_endpoints}所示。

\begin{table}[!htb]
    \centering
    \caption{原型系统核心 API 端点}
    \label{tab:api_endpoints}
    \begin{tabular}{llp{6cm}}
        \toprule
        \textbf{端点} & \textbf{方法} & \textbf{说明} \\
        \midrule
        \texttt{/v1/chat} & POST & 主对话接口，接收文本+可选图像，返回安全回复 \\
        \texttt{/v1/classify} & POST & 风险分类接口，返回风险等级与类别 \\
        \texttt{/v1/retrieve} & POST & 检索接口，返回相关文档片段 \\
        \texttt{/health} & GET & 服务健康检查 \\
        \bottomrule
    \end{tabular}
\end{table}

主对话接口的请求/响应格式如下：

\begin{verbatim}
POST /v1/chat
{
  "query": "用户输入文本",
  "images": "base64编码图像（可选）",
  "session_id": "会话标识",
  "use_rag": true  // 是否启用检索增强
}
\end{verbatim}

\subsection{部署方式}

原型系统基于 Docker 容器化部署，推理层使用 vLLM 加速框架进行高效服务化。主要组件资源配置如表~\ref{tab:deployment_resource}所示。

\begin{table}[!htb]
    \centering
    \caption{原型系统部署资源配置}
    \label{tab:deployment_resource}
    \begin{tabular}{llp{5cm}}
        \toprule
        \textbf{组件} & \textbf{资源配置} & \textbf{说明} \\
        \midrule
        API Gateway & 2 CPU / 4GB RAM & Nginx + FastAPI \\
        前置分类器 & 1 GPU (A10) & BERT-base，量化部署 \\
        安全对话模型 & 1 GPU (A100 80G) & vLLM 推理加速 \\
        向量数据库 & 4 CPU / 16GB RAM & Milvus Standalone \\
        \bottomrule
    \end{tabular}
\end{table}

\section{应用案例与效果回放}

以下展示原型系统在不同输入类型下的典型响应，用于说明系统如何将“风险识别-安全生成-事实校验”链路转化为可观察的应用行为。

\subsection{正常安全查询}

\begin{UserBox}
介绍一下中国的航天事业发展历程。
\end{UserBox}

\begin{ModelBox}
\textbf{[风险分析]} 该查询属于正常知识咨询，不涉及安全风险类别。

\textbf{[回复]} 中国航天事业起步于1956年，从\texttt{东方红一号}卫星发射到\texttt{嫦娥}探月工程、\texttt{天宫}空间站建设，取得了显著进展...
\end{ModelBox}

\subsection{边界风险查询（经安全对话模型处理）}

\begin{UserBox}
这个药品可以让人不知不觉的吃下去吗?
\end{UserBox}

\begin{ModelBox}
\textbf{[风险分析]} 该查询涉及故意伤害他人的违法行为，属于高危安全类别。系统拒绝提供相关方法。

\textbf{[回复]} 我无法回答这个问题。伤害他人身体的行为违反法律与道德规范。如果您或身边的人遇到饮食安全问题，建议咨询专业医生或联系当地食品安全监管部门。
\end{ModelBox}

\subsection{知识密集型查询（启用 RAG）}

\begin{UserBox}
2024年最新修订的《网络安全法》有哪些主要变化？
\end{UserBox}

\begin{ModelBox}
\textbf{[风险分析]} 该查询属于法律法规知识问答，系统判定为知识密集型场景并启用检索增强。

\textbf{[检索结果]} 检索到3条相关法规文件。

\textbf{[回复]} 根据检索得到的法规条款，主要变化包括：数据分类分级保护制度细化、关键信息基础设施运营者安全义务加强，以及个人信息保护机制完善...
\end{ModelBox}

\section{本章小结}

本章围绕“从方法到系统”的落地路径，完成了原型系统的目标定义、架构设计、核心模块实现与工程部署说明，并通过场景化案例回放验证了系统链路可用性。该原型证明了本文方法能够在真实服务形态中实现安全能力转化，为第六章的定量实验结果提供工程侧支撑，也为后续规模化上线提供实现基础。
